\begin{frame}{Постановка задач на практику}
    \begin{itemize}
        \item \textbf{Розробка архітектури:} створити модель на базі архітектури Mamba-SSM.% генеративну для оптимізації обчислень.
        \item \textbf{Розробка застосунку:} створити застосунок, де користувач обирає модель, дає фотографію та кольорові підказки і, як результат, отримує кілька варіантів колоризації.
        \item \textbf{Інтерактивність:} реалізувати механізм прийняття підказок з використанням Гауссівських розподілів для уникнення артефактів.
        \item \textbf{Багатоваріантність:} забезпечити здатність моделі генерувати декілька унікальних варіантів колоризації з одного набору вхідних даних.
        \item \textbf{Порівняльний аналіз:} провести оцінку рішення за системою об'єктивних метрик (PSNR, SSIM, LPIPS, FID, KID) та за ефективністю (VRAM, Time) відносно базових моделей (\href{https://arxiv.org/pdf/1705.02999}{SIGGRAPH17}, \href{https://arxiv.org/pdf/1603.08511}{ECCV16}, \href{https://arxiv.org/pdf/2212.11613}{DDColor}).
    \end{itemize}
\end{frame}